\documentclass[11pt,a4paper]{article}
% ══════════════════════════════════════════════════════════════════════════════
% Edexcel Further Mathematics (Core Pure 1) --- Shared LaTeX Preamble
% Elegant, Minimalist & Aesthetic Academic Design System
% ══════════════════════════════════════════════════════════════════════════════

\usepackage[a4paper, top=22mm, bottom=22mm, left=20mm, right=20mm]{geometry}
\usepackage{amsmath, amssymb, amsthm}
\usepackage{enumitem}
\usepackage{booktabs}
\usepackage{array}
\usepackage{parskip}
\usepackage{microtype}
\usepackage{fancyhdr}
\usepackage{titlesec}
\usepackage{mdframed}
\usepackage{xcolor}
\usepackage[hidelinks]{hyperref}
\usepackage{graphicx}
\usepackage{tikz}
\usepackage{pgfplots}
\pgfplotsset{compat=1.18}

% ── Refined Palette ───────────────────────────────────────────────────────────
\definecolor{primary}{gray}{0.10}
\definecolor{darkgray}{gray}{0.20}
\definecolor{medgray}{gray}{0.50}
\definecolor{lightgray}{gray}{0.96}
\definecolor{boxgray}{gray}{0.85}
\definecolor{accentline}{gray}{0.75}

% ── Section formatting ──────────────────────────────────────────────────────────
\titleformat{\section}[block]
  {\normalfont\bfseries\color{primary}}{}
  {0pt}{}[\vspace{1pt}{\color{accentline}\rule{\textwidth}{0.4pt}}\vspace{4pt}]
\titlespacing*{\section}{0pt}{14pt}{6pt}

% ── List formatting ─────────────────────────────────────────────────────────────
\setlist[enumerate,1]{label=\textbf{\arabic*.}, leftmargin=1.8em, itemsep=14pt}
\setlist[enumerate,2]{label=\textbf{(\alph*)}, leftmargin=1.8em, itemsep=6pt}

% ── Header / footer (Overlaps permanently eliminated) ─────────────────────────
\pagestyle{fancy}
\newcommand{\CPHeader}[2]{%
  \fancyhf{}%
  \setlength{\headheight}{22pt}%
  \renewcommand{\headrulewidth}{0.3pt}%
  \renewcommand{\headrule}{\hbox to\headwidth{\color{accentline}\leaders\hrule height \headrulewidth\hfill}}%
  \fancyhead[L]{\footnotesize\scshape\color{medgray} Core Pure 1: #1}%
  \fancyhead[R]{\footnotesize\color{medgray} Chapter \##2}%
  \renewcommand{\footrulewidth}{0.3pt}%
  \renewcommand{\footrule}{\hbox to\headwidth{\color{accentline}\leaders\hrule height \footrulewidth\hfill}}%
  \fancyfoot[L]{\footnotesize\color{medgray} Pearson Edexcel A-Level Further Mathematics $\cdot$ Core Pure 1}%
  \fancyfoot[C]{}% No center text to ensure zero collision with long titles
  \fancyfoot[R]{\footnotesize\color{medgray} \textbf{Henry.K} \enspace$\vert$\enspace Page \thepage}%
}

% ── Elegant Aesthetic Title Block ──────────────────────────────────────────────
\newcommand{\CPTitleBlock}[2]{%
  \vspace*{-8pt}
  \begin{center}
    {\footnotesize\scshape\color{medgray} Pearson Edexcel A-Level Further Mathematics \enspace$\cdot$\enspace Core Pure 1}\\[4pt]
    {\LARGE\bfseries\color{primary} #1}\\[4pt]
    {\small\color{darkgray} #2}\\[6pt]
    {\footnotesize\color{medgray} Problem Series authored by \textbf{Henry.K} \quad$\cdot$\quad 10 Exam Archetypes \quad$\cdot$\quad 65 Total Marks}\\[8pt]
    {\color{accentline}\rule{\textwidth}{0.5pt}}
  \end{center}
  \vspace{-6pt}
  \noindent{\footnotesize\itshape \textbf{Instructions:} Answer all ten questions. Full algebraic working must be shown for method marks; unsupported answers from calculators will not gain full credit. Give exact values in surd or fractional form unless specified otherwise.}\par
  \vspace{3pt}
  {\color{accentline}\rule{\textwidth}{0.5pt}}
  \vspace{8pt}
}

% ── Answer / Method / Investigate-further boxes ────────────────────────────────
\newmdenv[
  backgroundcolor=lightgray,
  linecolor=boxgray,
  linewidth=0.5pt,
  roundcorner=3pt,
  leftmargin=0pt, rightmargin=0pt,
  innerleftmargin=10pt, innerrightmargin=10pt,
  innertopmargin=6pt, innerbottommargin=6pt,
  skipabove=5pt, skipbelow=5pt
]{ansbox}

\newmdenv[
  backgroundcolor=lightgray!40,
  linecolor=darkgray,
  linewidth=1.2pt,
  leftline=true, rightline=false, topline=false, bottomline=false,
  leftmargin=0pt, rightmargin=0pt,
  innerleftmargin=10pt, innerrightmargin=8pt,
  innertopmargin=5pt, innerbottommargin=5pt,
  skipabove=5pt, skipbelow=7pt
]{invbox}

\newcommand{\ans}[1]{%
  \begin{ansbox}
    \textbf{Answer:} #1
  \end{ansbox}}

\newcommand{\method}[1]{%
  {\small\textbf{Method Criteria:}\par\vspace{2pt}#1}\par}

\newcommand{\inv}[1]{%
  \begin{invbox}
    \textbf{\small Deep Investigation:} {\small #1}
  \end{invbox}}

\newcommand{\qn}[1]{\textbf{#1.}\;}

\CPHeader{Series \& Summation}{3}

\begin{document}

\CPTitleBlock{Core Pure 1: Series \& Summation --- Answers \& Investigations}{Worked Solutions, Mark Criteria \& Deep Investigative Prompts}
\noindent\textit{Notice: Every question features a concise answer, detailed method criteria (M/A/B marks), and a dedicated investigation prompt to build structural insight.}

\section*{Section A \quad Procedural \& Direct Foundations}
\begin{enumerate}
%── Q1 [DIR] ─────────────────────────────────────────────────────────────
\item \textbf{[DIR]} \textit{Sum of the first $n$ positive integers with shifted lower limits}
\ans{[Final algebraic solution and exact values]}
\method{Key method step M1: apply definition or algebraic transformation. Accuracy mark A1: intermediate result. Final accuracy A1: exact form specified in question.}
\inv{Investigative exploration for Q1: What happens if the boundary condition or sign changes? How does this generalize to higher dimensions or arbitrary parameters?}

%── Q2 [DIR] ─────────────────────────────────────────────────────────────
\item \textbf{[DIR]} \textit{Sum of squares formula application and algebraic factorization}
\ans{[Final algebraic solution and exact values]}
\method{Key method step M1: apply definition or algebraic transformation. Accuracy mark A1: intermediate result. Final accuracy A1: exact form specified in question.}
\inv{Investigative exploration for Q2: What happens if the boundary condition or sign changes? How does this generalize to higher dimensions or arbitrary parameters?}

%── Q3 [DIR] ─────────────────────────────────────────────────────────────
\item \textbf{[DIR]} \textit{Sum of cubes formula and proving identities of the form $S_n = [f(n)]^2$}
\ans{[Final algebraic solution and exact values]}
\method{Key method step M1: apply definition or algebraic transformation. Accuracy mark A1: intermediate result. Final accuracy A1: exact form specified in question.}
\inv{Investigative exploration for Q3: What happens if the boundary condition or sign changes? How does this generalize to higher dimensions or arbitrary parameters?}

\end{enumerate}

\section*{Section B \quad Reverse Parameters \& Analysis}
\begin{enumerate}[resume]
%── Q4 [REV] ─────────────────────────────────────────────────────────────
\item \textbf{[REV]} \textit{Evaluating general cubic polynomials in $r$: sum of $(ar^3 + br^2 + cr + d)$}
\ans{[Final value of unknown parameter and corresponding equation]}
\method{Setup condition M1. Solving the resulting auxiliary polynomial or system M1. Accuracy criteria A1 with explicit rejection of inadmissible values.}
\inv{Investigative exploration for Q4: Check the geometric or physical interpretation when the discriminant is zero. Identify any bifurcation points.}

%── Q5 [REV] ─────────────────────────────────────────────────────────────
\item \textbf{[REV]} \textit{Parameter determination: finding constant $k$ given total sum $S_n$}
\ans{[Final value of unknown parameter and corresponding equation]}
\method{Setup condition M1. Solving the resulting auxiliary polynomial or system M1. Accuracy criteria A1 with explicit rejection of inadmissible values.}
\inv{Investigative exploration for Q5: Check the geometric or physical interpretation when the discriminant is zero. Identify any bifurcation points.}

%── Q6 [PRF] ─────────────────────────────────────────────────────────────
\item \textbf{[PRF]} \textit{Method of differences: standard partial fraction telescoping sum $\frac{1}{r(r+1)}$}
\ans{[Final value of unknown parameter and corresponding equation]}
\method{Setup condition M1. Solving the resulting auxiliary polynomial or system M1. Accuracy criteria A1 with explicit rejection of inadmissible values.}
\inv{Investigative exploration for Q6: Check the geometric or physical interpretation when the discriminant is zero. Identify any bifurcation points.}

\end{enumerate}

\section*{Section C \quad Proof, Geometry \& Multi-step}
\begin{enumerate}[resume]
%── Q7 [PRF] ─────────────────────────────────────────────────────────────
\item \textbf{[PRF]} \textit{Method of differences: stride-2 telescoping cancellation $\frac{1}{r(r+2)}$}
\ans{[Complete deductive proof or multi-part analytical solution]}
\method{Formal proof structure or systematic algebraic elimination. Explicit verification of base cases, cofactor signs, or orthogonal projections.}
\inv{Investigative exploration for Q7: Can this result be established via an alternative mathematical framework (e.g. geometric vs algebraic)?}

%── Q8 [PRF] ─────────────────────────────────────────────────────────────
\item \textbf{[PRF]} \textit{Method of differences with factorials or trigonometric identities}
\ans{[Complete deductive proof or multi-part analytical solution]}
\method{Formal proof structure or systematic algebraic elimination. Explicit verification of base cases, cofactor signs, or orthogonal projections.}
\inv{Investigative exploration for Q8: Can this result be established via an alternative mathematical framework (e.g. geometric vs algebraic)?}

%── Q9 [REV] ─────────────────────────────────────────────────────────────
\item \textbf{[REV]} \textit{Convergence and evaluation of infinite telescoping series}
\ans{[Complete deductive proof or multi-part analytical solution]}
\method{Formal proof structure or systematic algebraic elimination. Explicit verification of base cases, cofactor signs, or orthogonal projections.}
\inv{Investigative exploration for Q9: Can this result be established via an alternative mathematical framework (e.g. geometric vs algebraic)?}

%── Q10 [SYN] ─────────────────────────────────────────────────────────────
\item \textbf{[SYN]} \textit{Synoptic series applications linking summation to induction or polynomials}
\ans{[Complete deductive proof or multi-part analytical solution]}
\method{Formal proof structure or systematic algebraic elimination. Explicit verification of base cases, cofactor signs, or orthogonal projections.}
\inv{Investigative exploration for Q10: Can this result be established via an alternative mathematical framework (e.g. geometric vs algebraic)?}

\end{enumerate}

\end{document}